\subsection{Specialization with nonsemisimple residue algebra}
\label{subsec:closed-link}
\hypertarget{source-ordinary-closed-link}{}

One condition is still needed before the local prime comparison can be used along the
chain. Proposition~\ref{prop:marked-chain-support} requires that the specialized Cayley--Hamilton algebra remain nonsemisimple after extension of the residue field.
Let \(D\) be the ordinary determinant ring at the level, \(S\) its
ordered splitting algebra after the coefficient choices above and
\(B_0\to S\) the augmentation of the ring on the marked group.
Write \(G_0\) for the base CM group and \(G_*\twoheadrightarrow G_0\)
for the corresponding surjection from the marked CM group. Let \(P\)
be a prime of \(S\). Assume that over a field at the contracted
prime of \(D\), the specialization of the ordinary
Cayley--Hamilton algebra is nonsemisimple.
In the application its nonsemisimple quotient acts on the nonsplit
extension constructed in Section~\ref{sec:linked-components}.

\begin{proposition}[Nonsemisimplicity after specialization of the completed algebra]
\label{prop:marked-thick-realization}
Let \(P\subset S\) be a prime of the ordered splitting algebra
at which the ordinary Cayley--Hamilton algebra is nonsemisimple
over a field at the contraction to \(D\). For every prime \(Q\)
of a completed point ring \(B\) contracting to the preimage of
\(P\) under \(B_0\to S\), the specialization of the closed marked
relation quotient is nonsemisimple. In particular the mixed module satisfies 
\[
 \operatorname{Ann}_B M\subseteq Q.
\]

\end{proposition}

\begin{proof}
\proofstep{Nonsemisimplicity after scalar extension.}
First, for a finite-dimensional algebra over a field, a nonzero
radical is a nonzero nilpotent two-sided ideal. Extension of fields
preserves that ideal by faithful flatness. Thus nonsemisimplicity
persists over a common field. Two maps out of a ring with the same
prime kernel induce embeddings of the fraction field of its prime
quotient. The tensor product of the two target fields over that
fraction field is nonzero, and quotient by a prime followed by
passage to fractions gives a common field. If \(C=A/J\) presents the group algebra quotient before scalar extension to \(k'\), right exactness gives the specialized presentation
\[
 C\otimes k'\simeq (A\otimes k')/\operatorname{im}(J\otimes k').
\] These facts
give a nonsemisimple specialization over a field above \(P\). Finiteness here
follows from the spanning set for the integral Cayley--Hamilton
algebra in Section~\ref{sec:patching-matrices}.

\proofstep{Inclusion of the closed kernels before extending scalars.}
We need an inclusion of kernels \emph{before} the next field
extension. In \(S[G_0]\) form the ideal \(J'\) generated by the
image of the closed ordinary kernel from \(D[G_0]\). Put
\(C'=S[G_0]/J'\). The finite set of group elements spanning the
original closed quotient spans \(C'\) over \(S\). Continuity of the marked group algebra at each depth supplies an open normal subgroup with
\[
 U_b\triangleleft G_0,\qquad [u]-1\in\mathfrak m_S^b C'\quad(u\in U_b).
\]
In fact continuity on the finite original lattice gives this
assertion over \(D\), and \(D\to S\) is local. Also \(B_0\to S\)
is local, as can be seen on the common coefficient residue field.

Under augmentation, the determinant and its involution pull back from
the base. Applying the group and coefficient maps to a
Cayley--Hamilton relation thus gives the corresponding relation in
\(C'\). This includes polynomial-law relations with variable
coefficients: the base-extended law on the Cayley--Hamilton algebra
satisfies its characteristic polynomial after every scalar extension.
The local ordered products and the old inertia relations map to their
base relations, including their images under polarization. At an auxiliary place let \(\lambda_1,\dots,\lambda_n\)
be the ordered roots in \(S\) and \(\phi\) the image of
Frobenius. The image of the ordered relation is
\[
            \prod_{i=1}^{n}(\phi-\lambda_i)=0
\]
by Cayley--Hamilton for the actual characteristic
polynomial, keeping multiplicities. The tame generator
maps to \(1\) and the tame variables to zero, so the
image of each relation with a tame factor is also zero.
The involution gives the same for their polarized images.

\proofstep{The closed relation ideal and its specializations.}
Take now a member of the \emph{closure} of the marked relation ideal.
For each \(b\), the map to \(C'/\mathfrak m_S^bC'\)
kills the base image of \(U_b-1\) and the coefficient ideal
of depth \(b\). Therefore it factors through the group ring
on the corresponding finite quotient of \(G_*\) at that
depth. It kills each defining relation there by the preceding
calculation. For an element \(x\) of the closed ideal the induced map \(f\) therefore satisfies 
\[
 f(x)\bmod\mathfrak m_S^b C'=0\qquad(b\ge1).
\]
 The image lies in every power of the maximal ideal. Since \(C'\) is finite over the Noetherian local ring \(S\), Krull intersection gives
\[
 \bigcap_{b\ge1}\mathfrak m_S^b C'=0 .
\] This is the required inclusion of closed relation
kernels before extending the residue field.

Over the extension of \(\kappa(P)\) on which the base specialization is nonsemisimple, the surjection of groups induces
a surjection from the marked specialized group algebra modulo its
closed relations onto the base specialized algebra. Write \(E_{\rm mark}\twoheadrightarrow E_{\rm base}\) for the resulting quotient. The quotient of a semisimple finite-dimensional algebra is semisimple, so its contrapositive gives 
\[
 E_{\rm base}\text{ nonsemisimple}\quad\Longrightarrow\quad E_{\rm mark}\text{ nonsemisimple}.
\]
 

The common-field argument, this time over
\(B_0\) and with the fraction field at \(Q\), produces a field over which
the specialized algebra has a nonsemisimple quotient. The map from
\(B\) to this field has kernel \(Q\).
The completed group-algebra model of
Section~\ref{sec:patching-matrices} is the scalar extension of the
finite algebra with this closed presentation. It follows that
\(\operatorname{Ann}_B M\subseteq Q\) by the mixed-corner criterion
there. This proves the hypothesis on thick targets in the incidence
calculation, for links as well as the last target. The corresponding
statement for stable targets follows directly from the trace Gram
determinant after field extension. \label{cor:ordinary-closed-link-1}
\end{proof}
